% GENERATED FILE. Edit canonical lessons, then run npm run generate:books.
% canonical-source-hash: fnv1a64:237795098a70e00b
% canonical-lessons: SA-C509-jarjarah, SA-C509-slathah, SA-C509-randhram, SA-C509-syandate, SA-C509-jirnoddharah, SA-R509-review-502-509

\chapter{Something Needs Mending}
\label{ch:sa-something-needs-mending}
\hlchaptermodality{509}

\begin{chapteropening}
\textbf{By the end of this chapter:} \emph{I can say what is rickety, loose or dripping.}

Complete the last lesson of chapter 509: I can say what is rickety, loose or dripping.
\end{chapteropening}

\section[\texorpdfstring{\sk{जर्जरः} (jarjaraḥ)}{जर्जरः (jarjaraḥ)}]{\sk{जर्जरः} (jarjaraḥ) — rickety; worn to pieces}

\label{lesson:SA-C509-jarjarah}



\begin{quote}
Before the new one: say the Sanskrit for mends, then the Sanskrit for lime plaster.
\end{quote}

\subsection*{You'll want to know: \sk{जर्जरः}}
\textbf{\sk{जर्जरः}} — \emph{jarjaraḥ} — \textquotedblleft{}rickety; worn to pieces\textquotedblright{}.

\begin{quote}
\sk{मम} \sk{आसन्दी} \sk{जर्जरा} \sk{अस्ति।}
\end{quote}

\begin{quote}
\emph{mama āsandī jarjarā asti.}
\end{quote}

\begin{quote}
My chair is rickety.
\end{quote}

\begin{cousinweb}
\textbf{\sk{जर्जरः}} (\emph{jarjaraḥ}) is the adjective MW records as infirm, decrepit, decayed, broken in pieces; MW gives no derivation, so its origin is uncertain (MW).
\end{cousinweb}

\begin{grammarlens}[title={What you've built}]
One more word for describing what you lived and felt.
\end{grammarlens}

\subsection*{Your turn}
\begin{itemize}
\raggedright
  \item \emph{Say it:} \emph{jarjaraḥ}
  \item \emph{Say it:} \emph{jarjaraḥ}, once more
  \item \emph{Say it:} \emph{mama āsandī jarjarā asti.} — My chair is rickety.
  \item \emph{Recall:} say the Sanskrit for mends, then the Sanskrit for lime plaster, then say \emph{jarjaraḥ} again
\end{itemize}

\begin{tcolorbox}[breakable,colback=teal!4,colframe=teal!35!black,title={Before you move on}]
What does \sk{जर्जरः} mean? (\textquotedblleft{}Rickety; worn to pieces\textquotedblright{}.) Say it once more.
\end{tcolorbox}
\section[\texorpdfstring{\sk{श्लथः} (ślathaḥ)}{श्लथः (ślathaḥ)}]{\sk{श्लथः} (ślathaḥ) — slack; loose}

\label{lesson:SA-C509-slathah}



\begin{quote}
Before the new one: say the Sanskrit for lime plaster, then the Sanskrit for rickety.
\end{quote}

\subsection*{You'll want to know: \sk{श्लथः}}
\textbf{\sk{श्लथः}} — \emph{ślathaḥ} — \textquotedblleft{}slack; loose\textquotedblright{}.

\begin{quote}
\sk{द्वारस्य} \sk{कीलः} \sk{श्लथः} \sk{अस्ति।}
\end{quote}

\begin{quote}
\emph{dvārasya kīlaḥ ślathaḥ asti.}
\end{quote}

\begin{quote}
The nail in the door is loose.
\end{quote}

\begin{cousinweb}
\textbf{\sk{श्लथः}} (\emph{ślathaḥ}) comes from the root √ślath, \textquotedblleft{}to slacken, to become loose\textquotedblright{}; MW gives \emph{ślatha} as loose, relaxed, untied (MW).
\end{cousinweb}

\begin{grammarlens}[title={What you've built}]
One more word for describing what you lived and felt.
\end{grammarlens}

\subsection*{Your turn}
\begin{itemize}
\raggedright
  \item \emph{Say it:} \emph{ślathaḥ}
  \item \emph{Say it:} \emph{ślathaḥ}, once more
  \item \emph{Say it:} \emph{dvārasya kīlaḥ ślathaḥ asti.} — The nail in the door is loose.
  \item \emph{Recall:} say the Sanskrit for lime plaster, then the Sanskrit for rickety, then say \emph{ślathaḥ} again
\end{itemize}

\begin{tcolorbox}[breakable,colback=teal!4,colframe=teal!35!black,title={Before you move on}]
What does \sk{श्लथः} mean? (\textquotedblleft{}Slack; loose\textquotedblright{}.) Say it once more.
\end{tcolorbox}
\section[\texorpdfstring{\sk{रन्ध्रम्} (randhram)}{रन्ध्रम् (randhram)}]{\sk{रन्ध्रम्} (randhram) — an opening; a hole}

\label{lesson:SA-C509-randhram}



\begin{quote}
Before the new one: say the Sanskrit for rickety, then the Sanskrit for slack.
\end{quote}

\subsection*{You'll want to know: \sk{रन्ध्रम्}}
\textbf{\sk{रन्ध्रम्}} — \emph{randhram} — \textquotedblleft{}an opening; a hole\textquotedblright{}.

\begin{quote}
\sk{मम} \sk{स्यूते} \sk{एकं} \sk{रन्ध्रम्} \sk{अस्ति।}
\end{quote}

\begin{quote}
\emph{mama syūte ekaṃ randhram asti.}
\end{quote}

\begin{quote}
There is a hole in my bag.
\end{quote}

A neuter noun.

\begin{cousinweb}
\textbf{\sk{रन्ध्रम्}} (\emph{randhram}) probably comes from the root √rad, \textquotedblleft{}to split\textquotedblright{}, as MW suggests; MW gives \emph{randhra} as a slit, an opening, a hole (MW).
\end{cousinweb}

\begin{grammarlens}[title={What you've built}]
One more word for describing what you lived and felt.
\end{grammarlens}

\subsection*{Your turn}
\begin{itemize}
\raggedright
  \item \emph{Say it:} \emph{randhram}
  \item \emph{Say it:} \emph{randhram}, once more
  \item \emph{Say it:} \emph{mama syūte ekaṃ randhram asti.} — There is a hole in my bag.
  \item \emph{Recall:} say the Sanskrit for rickety, then the Sanskrit for slack, then say \emph{randhram} again
\end{itemize}

\begin{tcolorbox}[breakable,colback=teal!4,colframe=teal!35!black,title={Before you move on}]
What does \sk{रन्ध्रम्} mean? (\textquotedblleft{}An opening; a hole\textquotedblright{}.) Say it once more.
\end{tcolorbox}
\section[\texorpdfstring{\sk{स्यन्दते} (syandate)}{स्यन्दते (syandate)}]{\sk{स्यन्दते} (syandate) — trickles}

\label{lesson:SA-C509-syandate}



\begin{quote}
Before the new one: say the Sanskrit for slack, then the Sanskrit for an opening.
\end{quote}

\subsection*{You'll want to know: \sk{स्यन्दते}}
\textbf{\sk{स्यन्दते}} — \emph{syandate} — \textquotedblleft{}trickles\textquotedblright{}.

\begin{quote}
\sk{रन्ध्रात्} \sk{जलं} \sk{स्यन्दते।}
\end{quote}

\begin{quote}
\emph{randhrāt jalaṃ syandate.}
\end{quote}

\begin{quote}
Water trickles from the hole.
\end{quote}

Past (it …): \sk{जलं} \sk{स्यन्नवत्}, \emph{jalaṃ syannavat}.

\begin{cousinweb}
\textbf{\sk{स्यन्दते}} (\emph{syandate}) comes from the root √syand, \textquotedblleft{}to move on, to flow, to trickle\textquotedblright{}; its past participle is \emph{syanna} (MW).
\end{cousinweb}

\begin{grammarlens}[title={What you've built}]
One more word for describing what you lived and felt.
\end{grammarlens}

\subsection*{Your turn}
\begin{itemize}
\raggedright
  \item \emph{Say it:} \emph{syandate}
  \item \emph{Say it:} \emph{syandate}, once more
  \item \emph{Say it:} \emph{randhrāt jalaṃ syandate.} — Water trickles from the hole.
  \item \emph{Recall:} say the Sanskrit for slack, then the Sanskrit for an opening, then say \emph{syandate} again
\end{itemize}

\begin{tcolorbox}[breakable,colback=teal!4,colframe=teal!35!black,title={Before you move on}]
What does \sk{स्यन्दते} mean? (\textquotedblleft{}Trickles\textquotedblright{}.) Say it once more.
\end{tcolorbox}
\section[\texorpdfstring{\sk{जीर्णोद्धारः} (jīrṇoddhāraḥ)}{जीर्णोद्धारः (jīrṇoddhāraḥ)}]{\sk{जीर्णोद्धारः} (jīrṇoddhāraḥ) — a repair}

\label{lesson:SA-C509-jirnoddharah}



\begin{quote}
Before the new one: say the Sanskrit for an opening, then the Sanskrit for trickles.
\end{quote}

\subsection*{You'll want to know: \sk{जीर्णोद्धारः}}
\textbf{\sk{जीर्णोद्धारः}} — \emph{jīrṇoddhāraḥ} — \textquotedblleft{}a repair\textquotedblright{}.

\begin{quote}
\sk{एतस्य} \sk{गृहस्य} \sk{जीर्णोद्धारः} \sk{आवश्यकः} \sk{अस्ति।}
\end{quote}

\begin{quote}
\emph{etasya gṛhasya jīrṇoddhāraḥ āvaśyakaḥ asti.}
\end{quote}

\begin{quote}
This house needs repairs.
\end{quote}

A masculine noun.

\begin{cousinweb}
\textbf{\sk{जीर्णोद्धारः}} (\emph{jīrṇoddhāraḥ}) is built from \emph{jīrṇa}, \textquotedblleft{}worn out\textquotedblright{}, and \emph{uddhāra}, \textquotedblleft{}a raising up\textquotedblright{}; MW gives \emph{jīrṇoddhāra} as the repairing of what has fallen into ruin in a building (MW).
\end{cousinweb}

\begin{grammarlens}[title={What you've built}]
One more word for describing what you lived and felt.
\end{grammarlens}

\subsection*{Your turn}
\begin{itemize}
\raggedright
  \item \emph{Say it:} \emph{jīrṇoddhāraḥ}
  \item \emph{Say it:} \emph{jīrṇoddhāraḥ}, once more
  \item \emph{Say it:} \emph{etasya gṛhasya jīrṇoddhāraḥ āvaśyakaḥ asti.} — This house needs repairs.
  \item \emph{Recall:} say the Sanskrit for an opening, then the Sanskrit for trickles, then say \emph{jīrṇoddhāraḥ} again
\end{itemize}

\begin{tcolorbox}[breakable,colback=teal!4,colframe=teal!35!black,title={Before you move on}]
What does \sk{जीर्णोद्धारः} mean? (\textquotedblleft{}A repair\textquotedblright{}.) Say it once more.
\end{tcolorbox}
\section[(dialogue)]{Looking back: from In the Workshop to Something Needs Mending}

\label{lesson:SA-R509-review-502-509}



\begin{quote}
Say the Sanskrit for a workshop and a carpenter.
\end{quote}

\subsection*{Your turn}
Say these aloud:

\begin{itemize}
\raggedright
  \item a workshop, a carpenter, pares, a saw, a chisel
  \item a colleague, adroit, hard work, an employer, employs
  \item an application, a post, talent, lives by, skilled
  \item clay, earthen, kneads, a kiln, a shape
  \item weaves, a strand, a dyer, dyes, a tailor
  \item a bellows, tongs, melts down, blows on a fire, an ornament
  \item a brick, piles up, construction, mends, lime plaster
  \item rickety, slack, an opening, trickles, a repair
\end{itemize}

\begin{tcolorbox}[breakable,colback=teal!4,colframe=teal!35!black,title={Before you move on}]
A workshop? (\textbf{\sk{कर्मशाला}}, \emph{karmaśālā}.) A carpenter? (\textbf{\sk{तक्षकः}}, \emph{takṣakaḥ}.) Pares? (\textbf{\sk{तक्षति}}, \emph{takṣati}.) A saw? (\textbf{\sk{क्रकचः}}, \emph{krakacaḥ}.) A chisel? (\textbf{\sk{टंकः}}, \emph{ṭaṅkaḥ}.) A colleague? (\textbf{\sk{सहकारी}}, \emph{sahakārī}.) Adroit? (\textbf{\sk{दक्षः}}, \emph{dakṣaḥ}.) Hard work? (\textbf{\sk{परिश्रमः}}, \emph{pariśramaḥ}.) An employer? (\textbf{\sk{नियोक्ता}}, \emph{niyoktā}.) Employs? (\textbf{\sk{नियुंक्ते}}, \emph{niyuṅkte}.) An application? (\textbf{\sk{आवेदनम्}}, \emph{āvedanam}.) A post? (\textbf{\sk{नियुक्तिः}}, \emph{niyuktiḥ}.) Talent? (\textbf{\sk{प्रतिभा}}, \emph{pratibhā}.) Lives by? (\textbf{\sk{उपजीवति}}, \emph{upajīvati}.) Skilled? (\textbf{\sk{प्रवीणः}}, \emph{pravīṇaḥ}.) Clay? (\textbf{\sk{मृद्}}, \emph{mṛd}.) Earthen? (\textbf{\sk{मृन्मयः}}, \emph{mṛnmayaḥ}.) Kneads? (\textbf{\sk{मृद्नाति}}, \emph{mṛdnāti}.) A kiln? (\textbf{\sk{आपाकः}}, \emph{āpākaḥ}.) A shape? (\textbf{\sk{आकारः}}, \emph{ākāraḥ}.) Weaves? (\textbf{\sk{वयति}}, \emph{vayati}.) A strand? (\textbf{\sk{तन्तुः}}, \emph{tantuḥ}.) A dyer? (\textbf{\sk{रञ्जकः}}, \emph{rañjakaḥ}.) Dyes? (\textbf{\sk{रञ्जयति}}, \emph{rañjayati}.) A tailor? (\textbf{\sk{सौचिकः}}, \emph{saucikaḥ}.) A bellows? (\textbf{\sk{भस्त्रा}}, \emph{bhastrā}.) Tongs? (\textbf{\sk{सन्दंशः}}, \emph{sandaṃśaḥ}.) Melts down? (\textbf{\sk{विलापयति}}, \emph{vilāpayati}.) Blows on a fire? (\textbf{\sk{धमति}}, \emph{dhamati}.) An ornament? (\textbf{\sk{भूषणम्}}, \emph{bhūṣaṇam}.) A brick? (\textbf{\sk{इष्टका}}, \emph{iṣṭakā}.) Piles up? (\textbf{\sk{चिनोति}}, \emph{cinoti}.) Construction? (\textbf{\sk{निर्माणम्}}, \emph{nirmāṇam}.) Mends? (\textbf{\sk{सन्दधाति}}, \emph{sandadhāti}.) Lime plaster? (\textbf{\sk{सुधा}}, \emph{sudhā}.) Rickety? (\textbf{\sk{जर्जरः}}, \emph{jarjaraḥ}.) Slack? (\textbf{\sk{श्लथः}}, \emph{ślathaḥ}.) An opening? (\textbf{\sk{रन्ध्रम्}}, \emph{randhram}.) Trickles? (\textbf{\sk{स्यन्दते}}, \emph{syandate}.) A repair? (\textbf{\sk{जीर्णोद्धारः}}, \emph{jīrṇoddhāraḥ}.)
\end{tcolorbox}
